\chapter{Prototípus koncepciója}
\comment{A prototípus program célja annak demonstrálása, hogy a program elkészült, helyesen működik, valamennyi feladatát teljesíti. A prototípus változat egy elkészült program kivéve a kifejlett grafikus interfészt. Ez a program is parancssorból futtatható és karakteres ernyőkezelést alkalmaz. Az ütemezés, az aktív objektumok kezelése megoldott. A business objektumok -a megjelenítésre vonatkozó részeket kivéve -valamennyi metódusa a végleges algoritmusokat kell, hogy tartalmazza. A megjelenítés és működtetés egy alfanumerikus képernyőn vezérelhető és követhető, ugyanakkor a vezérlés fájlból is történhet és a megjelenítés fájlba is logolható, ezzel megteremtve a rendszer tesztelésének lehetőségét. Különös figyelmet kell fordítani a parancssori interfész logikájára, felépítésére, valamint arra, hogy az mennyiben tükrözi és teszi láthatóvá a program működését, a beavatkozások hatásait.}
\comment{\hrule}
\comment{Amennyiben változott a modell:}
\setcounter{section}{-1}
\pagebreak
\section{Változás hatása a modellre}
\subsection{Módosult osztálydiagram}
\comment{Az analízis modell osztálydiagramja a változások figyelembevételével.}
\diagram{docs/img/uml_class/zuzottko_class}{Módosult osztálydiagram}{18cm}
\subsection{Új vagy megváltozott metódusok}
\textbf{A bevezetett új osztályok}
\subsubsection{ZuzottkoSzoroFej}
\begin{class-template-responsibility}
    A rendkívül veszélyes jégpáncél csúszósságának a megszűntetésére szolgáló fej, amely a sávra szórja a fejben tárolt zúzottkövet. Fontos felelőssége, hogy a sávon lévő csapadékot nem távolítja el.
\end{class-template-responsibility}
\begin{class-template-baseclass}
    Fej \baseclass ZuzottkoSzoroFej
\end{class-template-baseclass}
\begin{class-template-attribute}
    \classitem{+fejAllapota: bool}{Jelzi, hogy az eszköz éppen be van-e kapcsolva.}
    \classitem{+zuzottkoKeszlet: int}{A fejben lévő kiszórható zúzottkő mennyisége.}
\end{class-template-attribute}
\begin{class-template-method}
    \classitem{+takaritHatas(s: Sav, h: Hokotro) : void (Override)}{Bekapcsolt állapotban kiszórja a zúzottkövet a sávra, ezáltal megakadályozva a jegen való megcsúszást. A sávon lévő csapadékot nem távolítja el.}
\end{class-template-method}

\textbf{A korábbi osztályokban történt módosítások}
\subsubsection{TakaritoJatekos}
\begin{class-template-responsibility}
    A hókotró flottát üzemeltető játékos reprezentációja. Szolgáltatásokat vesz igénybe a telephelyen (fejcsere, vásárlás, újratöltés), és karbantartja a gépeit.
\end{class-template-responsibility}
\begin{class-template-baseclass}
    Jatekos \baseclass TakaritoJatekos 
\end{class-template-baseclass}
\begin{class-template-method}
    \classitem{+hokotroVasarlas(Hokotro h) : bool}{Új hókotró vásárlása.}
    \classitem{+fejVasarlas(Fej f) : bool}{Új fej vásárlása.}
    \classitem{+fejLevetel(Hokotro h, Telephely t) : bool}{A hókotró jelenlegi fejének levétele és a raktárba helyezése a telephelyen.}
    \classitem{+fejFelteves(Hokotro h, Fej f) : bool}{Raktáron lévő fej felszerelése a hókotróra a telephelyen.}
    \classitem{+soToltes(Hokotro h) : int}{Pénzért feltölti a hókotró sókészletét a telephelyen.}
    \classitem{+kerozinToltes(Hokotro h) : int}{Pénzért feltölti a hókotró kerozinkészletét a telephelyen.}
    \classitem{+zuzottkoToltes(Hokotro h) : int}{Pénzért feltölti a hókotró zúzottkőkészletét a telephelyen.}
\end{class-template-method}

\subsection{Szekvencia-diagramok}
\diagram{docs/img/BMElogo}{Megváltozott szekvenciák}{3cm}
\diagram{docs/img/uml_sequence/zuzottkotoltes}{Zúzottkő töltése}{15cm}
\diagram{docs/img/uml_sequence/zuzottkofej}{Zúzottkő fej működése}{15cm}
\comment{\hrule}

\clearpage
\section{Prototípus interface-definíciója}
\comment{Definiálni kell a teszteket leíró nyelvet. Külön figyelmet kell fordítani arra, hogy ha a rendszer véletlen elemeket is tartalmaz, akkor a véletlenszerűség ki-bekapcsolható legyen, és a program determinisztikusan is tesztelhető legyen.}

\subsection{Az interfész általános leírása}
\comment{A protó (karakteres) input és output felületeit úgy kell kialakítani, hogy az input fájlból is vehető legyen illetőleg az output fájlba menthető legyen, vagyis kommunikációra csak a szabványos be- és kimenet használható.}

\subsection{Bemeneti nyelv}
\comment{A prototípus megvalósítása egy konzol alapú parancsértelmező rendszer keretein belül lesz megvalósítva. A parancsok kiadásán kívül lehetőség lesz azok hatásának átfogó követésére, hiszen a parancsok kimenetét és hatását azonnal látni fogja a felhasználó. A parancsok egyértelműen definiáltak lesznek, és a program képes lesz értelmezni azokat, valamint végrehajtani a megfelelő műveleteket. A parancsoknak lehetnek paramétereik. Magukról a parancsokról és paramétereikről a felhasználó ezen dokumentumból, illetve a help parancs kimenetelének tanulmányozásából szerezhet információt. A parancsok között szerepelnek olyanok, amelyek a játékosok menedzsmentjére, a hókotrók működtetésére, valamint a rendszer általános működésére vonatkoznak.}
\begin{itemize}
%INICIALIZÁLÓS PARANCSOK
    \item nehezseg [ertek]
    \begin{itemize}
        \item Leírás: Beállítja a játék nehézségét.
        \item \texttt{[érték]}: easy, medium, hard
    \end{itemize}

    \item letrehoz [tipus] [nev]
    \begin{itemize}
        \item Leírás: Létrehoz egy adott objektumot: típusa lehet TakarítóJátékos, BuszosJátékos, Hókotró (maximum egyet lehet létrehozni kezdésként), Busz, Autó
        \item \texttt{[tipus]}: TakarítóJátékos, BuszosJátékos, Hókotró, Busz, Autó
        \item \texttt{[nev]}: egyedi azonosító, amely a későbbiekben a parancsokban használható
    \end{itemize}

    \item szerkeszt [nev] [ujnev] 
    \begin{itemize}
        \item Leírás: Változtatja a korábban létrehozott objektum nevét.
        \item \texttt{[nev]}: a szerkesztendő objektum neve
        \item \texttt{[ujnev]}: az új név, amely egyedi azonosító kell, hogy legyen
    \end{itemize}

    \item kivalaszt [nev]
    \begin{itemize}
        \item Leírás: Kiválasztja a megadott nevű objektumot, amelyre a további parancsok hatni fognak.
        \item \texttt{[nev]}: A kiválasztandó objektum neve
    \end{itemize}

    \item betolt [fajlnev]
    \begin{itemize}
        \item Leírás: Betölti a megadott fájlból az objektumokat.
        \item \texttt{[fajlnev]}: A betöltendő fájl neve
    \end{itemize}

    %MOZGATÓ PARANCSOK

    \item lepes [ut] [sav]
    \begin{itemize}
        \item Leírás: Egy jármű mozgatása a pályán, egy szomszédos útra. 
        \item \texttt{[ut]} : Az út neve amire szeretne lépni a játékos, amelynek a jelenlegi út szomszédságában kell lennie.
        \item \texttt{[sav]} : Opcionális paraméter, amely megadja, hogy a jármű melyik sávba szeretne lépni az adott úton, amennyiben az út több sávos.
    \end{itemize}

    \item wait [nev]
    \begin{itemize}
        \item Leírás: Lépés kihagyása, a játékos nem hajt végre lépést.
    \end{itemize}

    %Menedzsment parancsok

    \item hokotrovasarlas [nev]
    \begin{itemize}
        \item Leírás: A játékos megvásárol egy új hókotró járművet, amennyiben telephelyen tartózkodik és rendelkezik a megfelelő összeggel.
        \item \texttt{[nev]}: a megvásárolni kívánt hókotró neve, amely egyedi azonosító kell, hogy legyen
    \end{itemize}

    \item allapot [nev]
    \begin{itemize}
        \item Leírás: Kiírja egy adott objektum (játékos, hókotró, út) aktuális tulajdonságait.
        \item \texttt{[nev]}: az állapot lekérdezni kívánt objektum neve
    \end{itemize}

    \item lista [tipus]
    \begin{itemize}
        \item Leírás: Kiírja egy adott objektum (játékos, hókotró, út) aktuális tulajdonságait.
        \item \texttt{[tipus]}: az állapot lekérdezni kívánt objektum neve
    \end{itemize}
    
    \item terkep
    \begin{itemize}
        \item Leírás: Szövegesen kilistázza a csomópontokat és a köztük lévő utakat, hogy lehessen tudni, honnan hova lehet lépni.
    \end{itemize}
    
%HÓKOTRÓS PARANCSOK
    \item fejcsere [hokotro] [ujfej]
    \begin{itemize}
        \item Leírás: Kicseréli a hókotró fejét egy új fejre, amennyiben a játékos rendelkezik a megfelelő fejjel a készletében, és telephelyen tartózkodik.
        \item \texttt{[hokotro]}: a jelenlegi hókotró, amin fejet szeretnénk cserélni
        \item \texttt{[ujfej]}: a kívánt új fej típusa, amire cserélni szeretnénk a jelenlegi fejet
    \end{itemize}

    \item fejvasarlas [hokotro] [fejtipus]
    \begin{itemize}
        \item Leírás: Új fejet vásárol a játékos a készletére, amennyiben telephelyen tartózkodik és rendelkezik a megfelelő összeggel.
        \item \texttt{[hokotro]}: a jelenlegi hókotró, amelyhez fejvásárlást szeretnénk végrehajtani
        \item \texttt{[fejtipus]}: a kívánt fej típusa, amelyet vásárolni szeretnénk
    \end{itemize}

    \item sotoltes [hokotro]
    \begin{itemize}
        \item Leírás: Amennyiben a játékos telephelyen tartózkodik, feltölti a hókotró sókészletét maxra. 
        \item \texttt{[hokotro]}: a jelenlegi hókotró, amelynek sókészletét szeretnénk feltölteni
    \end{itemize}

    \item kerozintoltes [hokotro]
    \begin{itemize}
        \item Leírás: Amennyiben a játékos telephelyen tartózkodik, feltölti a hókotró kerozinkészletét maxra.
        \item \texttt{[hokotro]}: a jelenlegi hókotró, amelynek kerozinkészletét szeretnénk feltölteni
    \end{itemize}

    \item zuzalektoltes [hokotro]
    \begin{itemize}
        \item Leírás: Amennyiben a játékos telephelyen tartózkodik, feltölti a hókotró zuzalékkészletét maxra.
        \item \texttt{[hokotro]}: a jelenlegi hókotró, amelynek zuzalékkészletét szeretnénk feltölteni
    \end{itemize}

    \item takarit [hokotro] [sav]
    \begin{itemize}
        \item Leírás: A hókotro vévigmegy két csomópont között, és takarít az adott sávon. 
        \item \texttt{[hokotro]}: a jelenlegi hókotró, amelyik takarítani fog
        \item \texttt{[sav]}: az út azon sávja, amit takarítani tervezünk
    \end{itemize}

%RENDSZER PARANCSOK
    \item havazas 
    \begin{itemize}
        \item Leírás: Havazás eseményét indítja el, amelynek hatására a pályán hó jelenik meg. 
    \end{itemize}

    \item kilepes
    \begin{itemize}
        \item Leírás: Kilép a játékból
    \end{itemize}

    \item help
    \begin{itemize}
        \item Leírás: Megjeleníti a parancsok listáját
    \end{itemize}


\end{itemize}
Példák a bemeneti nyelvre:
\begin{verbatim}
    nehezseg medium
    betolt palya_alap.txt
    terkep
    lista Hokotro
    lista Auto

    letrehoz TakaritoJatekos takarito1
    letrehoz BuszosJatekos buszsofor1
    letrehoz Hokotro alap_kotro
    letrehoz Busz bkv_busz
    szerkeszt alap_kotro kotro_01
    allapot kotro_01
\end{verbatim}

\subsection{Kimeneti nyelv}
\comment{Egyértelműen definiálni kell, hogy az egyes bemeneti parancsok végrehajtása után előálló állapot milyen formában jelenik meg a szabványos kimeneten. A program képes legyen olyan kimenetet előállítani, amellyel az objektumok állapota ellenőrizhető (pl. save). Ebben az alfejezetben is precízen definiálni kell, hogy a kimenet nyelve milyen elemekből és milyen szintakszissal áll elő.}

A kimeneti üzenetek négy fő kategóriába sorolhatók:

\begin{itemize}
    \item \textbf{Visszajelzések (Siker / Hiba):} Minden kiadott parancs azonnali visszajelzést generál a végrehajtás eredményéről.
    \begin{itemize}
        \item \textbf{Szintaxis (Siker):} \texttt{> OK: [Rövid leírás a végrehajtott műveletről]}
        \item \textbf{Szintaxis (Hiba):} \texttt{> ERROR: [A hiba pontos oka, pl. "Nincs elegendő pénz" vagy "Nincs telephelyen"]}
    \end{itemize}

    \item \textbf{Állapot lekérdezések (allapot, lista):} Determinisztikus, kulcs-érték párokat tartalmazó reprezentáció, amely egyaránt alkalmas a program aktuális állapotának fájlba mentésére (save) és annak későbbi visszaellenőrzésére.
    \begin{itemize}
        \item \textbf{Játékos szintaxis:} \texttt{[Típus] [Név] | Pénz:[Összeg] | Járművek:[jármű1, jármű2, ...] | Raktár:[fej1, fej2, ...]}
        \item \textbf{Hókotró szintaxis:} \texttt{Hokotro [Név] | Poz:[Út neve, Sáv indexe] | Fej:[Aktív fej típusa] | Készletek:[So: X, Kerozin: Y, Zuzottko: Z] | Állapot:[Aktív/Baleset(hátralévő körök)]}
        \item \textbf{Út és Sáv szintaxis:} \texttt{Ut [Név] ([Típus]) | Szomszédok:[csomópont1, csomópont2] | Sávok: [Sáv0(Hó: X, Jég: Y, Sózott: I/N)], [Sáv1(Hó: X, Jég: Y, Sózott: I/N)]}
    \end{itemize}

    \item \textbf{Térkép vizualizáció (terkep):} Hierarchikus listázás, amely megmutatja a csomópontokat és az azokat összekötő utakat.
    \begin{itemize}
        \item \textbf{Szintaxis:} 
        \begin{verbatim}
Csomópont [Név] ([Típus])
  -> [Szomszédos Csomópont Név] : [Út Név] ([Sávok száma] sáv)
        \end{verbatim}
    \end{itemize}

    \item \textbf{Események és logok (havazas, baleset, kör vége):} A rendszer működése során automatikusan bekövetkező események (pl. NPC lépések, megcsúszás) logolása.
    \begin{itemize}
        \item \textbf{Szintaxis:} \texttt{[ESEMÉNY] [Esemény pontos leírása és az érintett objektumok]}
    \end{itemize}
\end{itemize}

\textbf{Példák a kimeneti nyelvre a bemenetek alapján:}

\begin{verbatim}
// Bemenet: letrehoz TakaritoJatekos takarito1
> OK: TakaritoJatekos 'takarito1' sikeresen letrehozva.

// Bemenet: allapot kotro_01
Hokotro kotro_01 | Poz: Fo_ut, 0 | Fej: SoproFej | Készletek: So:0, Kerozin:0, Zuzottko:0 | Állapot: Aktív

// Bemenet: fejvasarlas kotro_01 SarkanyFej
> ERROR: A jármű nem a telephelyen tartózkodik, a vásárlás meghiúsult.

// Bemenet: kerozintoltes kotro_01
> OK: 'kotro_01' kerozinkészlete feltöltve (Aktuális: 100%). Pénz levonva: 50.

// Bemenet: takarit kotro_01 0
> OK: 'kotro_01' végighaladt a 'Fo_ut' 0. sávján.
[ESEMÉNY] SarkanyFej bekapcsolva. Kerozin szint csökkent: 10%.
[ESEMÉNY] Fo_ut 0. sávján 2 egység hó elolvasztva.

// Bemenet: havazas
> OK: Havazás szimuláció elindítva.
[ESEMÉNY] Fo_ut összes sávján a hóréteg +1 egységgel nőtt.
[ESEMÉNY] Alagut_1 sávjain a hóréteg nem változott.

// Bemenet: terkep
Csomópont Telephely_1 (Telephely)
  -> Csomópont_A : Fo_ut (2 sáv)
Csomópont_A (Kereszteződés)
  -> Telephely_1 : Fo_ut (2 sáv)
  -> Csomópont_B : Mellek_ut (1 sáv)
\end{verbatim}

\section{Összes részletes use-case}

\begin{use-case}
    {Biokerozin töltés}
    {Biokerozin töltés a Sárkány fejbe a TakarítóManager által kezdeményezve}
    {TakarítóManager}
    Hasonlóan a sóhoz, a sárkány fej működtetéséhez szükséges biokerozin is a Telephelyen vételezhető. A takarító szerepet játszó játékos pénzért cserébe tölti fel a munkagép üzemanyagtartályát, biztosítva a gázturbina további működését. A rendszer a tranzakció sikere után frissíti a hókotró biokerozin-készletét, lehetővé téve a takarítás végrehajtását.
\end{use-case}

\begin{use-case}
    {Busz haladása és fizetés}
    {Busz haladása és fizetés a BuszsoforManager által kezdeményezve}
    {BuszsoforManager}
    A BuszsoforManager szerepkörben tevékenykedő játékos kezdeményezi a busz léptetését az úthálózaton. A jármű a megadott cél felé halad, miközben a rendszer ellenőrzi, hogy az adott kereszteződés egy Végállomás-e. Amennyiben a busz végállomásra ér, a program regisztrálja a megtett kört, növeli a teljesített utak számát, és a játékos vagyona gyarapodik a menetdíjból származó bevétellel.
\end{use-case}

\begin{use-case}
    {Hókotró fejcsere}
    {Hókotró fejének cseréje a TakarítóManager által kezdeményezve}
    {TakarítóManager}
    A TakarítoManager a Telephelyen kezdeményezi a Hókotróra szerelt munkaeszköz cseréjét. A folyamat során a rendszer ellenőrzi, hogy a Jármű a telephelyen tartózkodik-e, majd a régi fejet leszerelik és a raktárba helyezik. Ezt követően, ha a Játékos rendelkezik a felszerelni kívánt új fejjel, az rákerül a gépre, és törlődik a raktárkészletből, hiszen onnantól az lesz az aktív fej a hókotrón.
\end{use-case}

\begin{use-case}
    {Fej vásárlása}
    {Új hókotró fej vásárlása a TakarítóManager által kezdeményezve}
    {TakarítóManager}
    A TakarítoManager a Telephelyen speciális takarítóeszközöket (például Sárkány- vagy Sószóró fejet) vásárolhat a hatékonyabb munkavégzés érdekében. A vásárlás során a rendszer ellenőrzi a pénzügyi keretet, majd a vételár kifizetése után az új eszköz bekerül a játékos raktárába. A megvásárolt fejeket a játékos később bármelyik Hókotrójára felszerelheti a telephelyen.
\end{use-case}

\begin{use-case}
    {Hányó fej működése}
    {Hányó fej működése a jármű mozgatása közben}
    {JátékosManager}
    A fej először ellenőrzi, hogy található-e hóréteg vagy feltört jég az úton. Amennyiben van eltávolítandó csapadék, a fej azt nagy erővel az út mellé szórja, így a hó véglegesen kikerül a rendszerből, és nem kerül át szomszédos sávokba. A tömör jégpáncél viszont változatlanul marad a sávon. 
\end{use-case}

\begin{use-case}
    {Havazás szimulációja}
    {Havazás szimulációja minden pillanatban}
    {Tesztelő}
    A környezeti hatások modellezéséért felelős Úthálózat minden körben végrehajtja a havazás folyamatát. A rendszer egyenletesen növeli a hóvastagságot a város minden útszakaszán, kivéve az alagutakat, amelyeket a kialakításuk megvéd a csapadéktól. 
\end{use-case}

\begin{use-case}
    {Játék vége 1 – Elakadt tömegközlekedés}
    {Játék vége 1 – Elakadt tömegközlekedés a városban}
    {Tesztelő}
    A rendszer minden lépés után ellenőrzi a forgalom állapotát és a játék folytathatóságát. Ez a use case akkor következik be, ha a rendszer megállapítja, hogy a városban egyetlen busz sem képes a mozgásra (például balesetek vagy hótorlaszok miatt). Mivel a tömegközlekedés megbénulása a város működésképtelenségét jelenti, a játék ebben a pillanatban vereséggel véget ér.
\end{use-case}

\begin{use-case}
    {Játék vége 2 – Kijárási tilalom}
    {Játék vége 2 – Kijárási tilalom a városban}
    {Tesztelő}
    A rendszer szintén folyamatosan monitorozza a városban bekövetkezett balesetek és úttorlaszok számát. Amennyiben az összecsúszott és mozgásképtelenné vált járművek száma elér egy előre meghatározott kritikus határértéket, “kijárási tilalmat hirdetnek”. Ez a vészhelyzet a szimuláció azonnali leállítását és a játék végét vonja maga után.
\end{use-case}

\begin{use-case}
    {Autó áthaladása és jéggé tömörülés}
    {Autó áthaladása és jéggé tömörülés a sávon}
    {Tesztelő}
    Amikor egy NPC autó vagy egy busz áthalad egy havas útsávon, súlyával tömöríti az ott lévő havat. A sáv minden áthaladáskor növeli a számlálóját, és amint eléri a kritikus (5 jármű) határt, a sávon lévő teljes hómennyiség azonnal jégpáncéllá tömörödik. Ez a folyamat dinamikusan változtatja az útviszonyokat, növelve a balesetveszélyt a forgalom számára.
\end{use-case}

\begin{use-case}
    {Jégtörő fej működése}
    {Jégtörő fej működése a sávon lévő jégpáncél feltörésére}
    {JátékosManager}
    A fej feladata a sávon lévő masszív jégpáncél fizikai feltörése, amelyet “feltört jéggé” (a modellben hóvá) alakít át, de az útról nem távolítja el, így az visszatömörödhet jéggé. 
\end{use-case}

\begin{use-case}
    {Leszórt só}
    {Leszórt só hatása a sáv állapotára}
    {Tesztelő}
    A sószóró fej által az útra juttatott só aktív hatást fejt ki a sáv állapotára. Felolvasztja a havat és a jeget, illetve meggátolja azok újboli kialakulását egy bizonyos ideig. 
\end{use-case}

\begin{use-case}
    {Jármű megcsúszása}
    {Jármű megcsúszása a jeges sávon}
    {Tesztelő}
    Amikor egy jármű (autó vagy busz) jeges sávra lép, a rendszer baleseti kalkulációt végez a globális nehézségi szint és a jégvastagság alapján. Sikertelen kimenetel esetén a jármű megcsúszik, elszenvedi a balesetet és meghatározott ideig mozgásképtelenné válik. Az ilyen jármű eltorlaszolja az adott sávot, akadályozva a forgalom többi résztvevőjét, amíg a mozgásképtelensége le nem jár.
\end{use-case}

\begin{use-case}
    {Sárkány fej működése}
    {Sárkány fej működése a sávon lévő hó és jég eltávolítására}
    {JátékosManager}
    A fej működésének feltétele, hogy a Sárkányfej bekapcsolt állapotban legyen és rendelkezzen elegendő biokerozinnal. Aktív állapotban a gép biokerozint fogyaszt, cserébe az adott sávon minden havat és jégpáncélt azonnal és véglegesen elolvaszt, megtisztítva az utat a forgalom számára.
\end{use-case}

\begin{use-case}
    {Söprő fej működése}
    {Söprő fej működése a sávon lévő hó eltávolítására}
    {JátékosManager}
    A fej lekérdezi az aktuális sáv hóvastagságát, majd a havat és a feltört jeget a menetirány szerinti jobb oldali szomszédos sávba továbbítja. Ha a gép a legszélső sávban halad, a csapadék az út mellé kerül, ahol már nem akadályozza tovább a közlekedést.
\end{use-case}

\begin{use-case}
    {Sószóró fej működése}
    {Sószóró fej működése a sávon lévő hó és jég olvasztására és újraképződésének megakadályozására}
    {JátékosManager}
    Ez a fej működés közben csökkenti a gép sókészletét, és a sávhoz egy speciális só-objektumot rendel. Ez az anyag nemcsak felolvasztja a meglévő havat és jeget, hanem pár körön keresztül megakadályozza az újabb csapadék megmaradását is az adott sávon. 
\end{use-case}

\begin{use-case}
    {Só töltése a sószóró fejbe}
    {Só töltése a sószóró fejbe a TakarítóManager által kezdeményezve}
    {TakarítóManager}
    A hókotrókat vezető játékos a Telephelyen kezdeményezi a hókotró sószóró fejéhez tartozó sókészlet feltöltését. A folyamat során a rendszer ellenőrzi a játékos anyagi fedezetét, majd a megfelelő összeg levonása után a jármű sókészlete maximális szintre áll vissza. Ez a művelet elengedhetetlen a sószóró fej folyamatos és hatékony üzemeltetéséhez a jeges útszakaszokon.
\end{use-case}

\begin{use-case}
    {Új hókotró vásárlása}
    {Új hókotró vásárlása a TakarítóManager által kezdeményezve}
    {TakarítóManager}
    A Játékos a Telephelyen új Hókotrót vásárolhat a vagyona terhére, hogy bővítse a flottáját. A rendszer ellenőrzi a tartózkodási helyet és az anyagi fedezetet. Ha minden feltétel teljesül, levonja a vételárat a játékos pénzéből. Sikeres tranzakció után létrejön az új hókotró objektum, amely azonnal munkára fogható.
\end{use-case}

\section{Tesztelési terv}
\comment{A tesztelési tervben definiálni kell, hogy a be- és kimeneti fájlok egybevetésével miként végezhető el a program tesztelése. Meg kell adni teszt forgatókönyveket. Az egyes teszteket elég informálisan, szabad szövegként leírni. Teszt-esetenként egy-öt mondatban. Minden teszthez meg kell adni, hogy mi a célja, a proto mely funkcionalitását, osztályait stb. teszteli. Az alábbi táblázat minden teszt-esethez külön-külön elkészítendő.}

% ==========================================
% A tesztelés megvalósítása és a be-/kimeneti fájlok egybevetése
% ==========================================
% A prototípus tesztelése determinisztikus forgatókönyvek alapján történik. 
% Minden tesztesethez tartozik egy előre megírt bemeneti parancsfájl (input.txt), 
% amely a konzolos menüpontok kiválasztását és a program által feltett kérdésekre 
% adott válaszokat tartalmazza. A program futtatása során ezt a fájlt a 
% --input=input.txt paraméterrel adjuk át. A futás során keletkező konzolos üzeneteket, 
% logolásokat és állapotváltozásokat a --output=actual_output.txt paraméterrel fájlba irányítjuk. 
% A tesztelés kiértékelése a létrejött kimeneti fájl és egy előzetesen, manuálisan verifikált 
% elvárt kimeneti fájl (expected_output.txt) soronkénti összehasonlításával (diff / fc) történik. 
% Ha a két fájl karaktermillióra megegyezik, a teszt sikeres.

% ==========================================
% 1. Játékosok menedzsment akciói (Vásárlás, felszerelés, töltés)
% ==========================================

\test{Hókotró fejének cseréje - Sikeres}
{A Takarító játékos a telephelyen tartózkodva lecseréli a hókotrója jelenlegi fejét egy másikra, amely elérhető a raktárban.}
{A fejcsere metódusainak (\texttt{TakaritoManager}, \texttt{Hokotro}) helyes lefutását és a régi, valamint új \texttt{Fej} objektumkapcsolatok frissülését teszteli.}

\test{Hókotró fejének cseréje - Sikertelen (Nincs telephelyen)}
{A játékos a város egy normál útszakaszán próbál meg fejet cserélni a hókotróján.}
{Annak ellenőrzése, hogy a rendszer elutasítja a cserét, mivel a művelet kizárólag a \texttt{Telephely} csomóponton hajtható végre.}

\test{Új hókotró vásárlása - Sikeres}
{A Takarító menedzser a telephelyen állva új hókotrót vásárol, és rendelkezik a vásárláshoz szükséges elegendő pénzzel.}
{A vásárlási feltételek ellenőrzése, a pénzlevonás, és a \texttt{TakaritoManager} járműlistájának bővülésének validálása.}

\test{Új hókotró vásárlása - Sikertelen (Nincs elég pénz)}
{A menedzser a telephelyen próbál hókotrót vásárolni, de a kasszában nincs meg a fedezet.}
{Validálja, hogy a vagyon nem megy mínuszba, és a flotta nem bővül járművel fedezet hiányában.}

\test{Új hókotró vásárlása - Sikertelen (Nincs telephelyen)}
{A menedzser egy útszakaszon tartózkodik, és elegendő pénz birtokában próbál hókotrót vásárolni.}
{Ellenőrzi, hogy a helyszíni feltétel hiányában a vásárlás meghiúsul.}

\test{Új fej vásárlása - Sikeres}
{A menedzser egy speciális takarítófejet (pl. Zúzottkőszóró fej) vásárol a telephelyen, elegendő pénzből.}
{A fejvásárlás menetének, a készletkezelésnek (a fej a raktárba kerül) és a pénzkövetésnek az ellenőrzése.}

\test{Új fej vásárlása - Sikertelen (Nincs elég pénz)}
{A menedzser a telephelyen próbál meg fejet vásárolni, de nincs rá elég pénze.}
{A tranzakció sikertelen lefutásának ellenőrzése, a raktárkészlet változatlanságának validálása.}

\test{Új fej vásárlása - Sikertelen (Nincs telephelyen)}
{A menedzser nem telephelyen állva próbál takarítófejet venni a raktárába.}
{A tranzakció elutasításának ellenőrzése a helyszín miatt.}

\test{Fogyóeszköz töltése - Só (Sikeres)}
{A játékos a telephelyen sót tölt a hókotróra szerelt sószóró fejbe.}
{A sótöltés logikájának, a készletváltozásnak és a \texttt{SoszoroFej} belső állapotfrissítésének ellenőrzése.}

\test{Fogyóeszköz töltése - Só (Sikertelen, nincs telephelyen)}
{A játékos útközben, a várost járva próbál sót tölteni a sószóró fejbe.}
{Az akció blokkolásának ellenőrzése helyszíni megkötés miatt.}

\test{Fogyóeszköz töltése - Biokerozin (Sikeres)}
{A játékos a telephelyen biokerozint tölt a sárkány fejbe.}
{A kerozin feltöltésének és a \texttt{SarkanyFej} kapacitásnövekedésének validálása.}

\test{Fogyóeszköz töltése - Biokerozin (Sikertelen, nincs telephelyen)}
{A játékos munka közben próbál biokerozint vételezni a sárkány fejbe.}
{A művelet meghiúsulásának ellenőrzése.}

\test{Fogyóeszköz töltése - Zúzottkő (Sikeres)}
{A játékos a telephelyen zúzottkövet tölt a zúzottkőszóró fejbe.}
{A \texttt{ZuzottkoszoroFej} belső készletének sikeres feltöltésének tesztelése.}

\test{Fogyóeszköz töltése - Zúzottkő (Sikertelen, nincs telephelyen)}
{A játékos útszakaszon próbálja feltölteni a zúzottkőszóró fejét.}
{A művelet meghiúsulásának validálása.}

% ==========================================
% 2. Tömegközlekedés és forgalom mozgása
% ==========================================

\test{Busz mozgása - Végállomás elérése (Sikeres kör)}
{A Buszsofőr egy buszt egy olyan csomópontra léptet, amely a busz útvonalának végállomása.}
{A \texttt{Busz} és a \texttt{BuszsoforManager} tesztelése: ellenőrzi a körtelesítés regisztrációját és a bevétel jóváírását.}

\test{Busz mozgása - Köztes csomópont (Nincs körtelesítés)}
{A buszsofőr egy buszt mozgat, de a cél egy normál útszakasz vagy köztes csomópont.}
{Validálja, hogy köztes lépés esetén a rendszer nem regisztrál teljesített kört és nem ad pénzjutalmat.}

\test{Jármű elakadása magas hóban (Sikertelen haladás)}
{Egy jármű (Busz vagy NPC Autó) egy olyan sávra lépne, ahol a hó vastagsága mozgást akadályozó szintű.}
{A normál \texttt{Jarmu} osztályok lépésének blokkolásának tesztelése a \texttt{Sav} magas csapadékszintje miatt.}

\test{Hókotró mozgatása (Sikeres haladás extrém időben is)}
{A Takarító játékos hókotróját egy vastag hóréteggel vagy jégpáncéllal borított sávra irányítja.}
{A \texttt{Hokotro} környezeti hatásokkal (elakadás, megcsúszás) szembeni immunitásának ellenőrzése: a jármű akadálytalanul halad tovább.}

\test{Jármű megcsúszása jeges úton (Baleset)}
{Egy normál jármű jégpáncélos sávon halad át, ahol a baleseti ráta alapján megcsúszás következik be.}
{A balesetkalkuláció és az ütközési logika tesztelése: a jármű baleseti állapotba kerülésének validálása.}

\test{Autó áthaladása és jéggé tömörülés}
{Egy önműködő NPC autó áthalad egy havas sávon.}
{A \texttt{Sav.autoAthalad()} metódus tesztelése: validálja, hogy a hó a jármű súlya alatt \texttt{Jegpancel}-lá tömörül.}

% ==========================================
% 3. Takarítófejek működése
% ==========================================

\test{Sárkány fej működése - Sikeres olvasztás}
{A sárkány fejjel felszerelt hókotró áthalad egy havas/jeges sávon, és van elég biokerozinja.}
{Az azonnali hó- és jégolvasztó hatás, valamint a kerozinkészlet csökkenésének ellenőrzése.}

\test{Sárkány fej működése - Sikertelen (Kifogyott kerozin)}
{A sárkány fej áthalad a csapadékos sávon, de az üzemanyagtartálya üres.}
{Ellenőrzi, hogy a fej kerozin hiányában nem olvasztja el a csapadékot, a sáv állapota változatlan marad.}

\test{Sószóró fej működése - Sikeres szórás}
{A sószóró fej áthalad a sávon, van elég sója, így kiszórja azt.}
{A só kihelyezésének, a meglévő hó olvasztásának és az útburkolat ideiglenes tisztántartásának (\texttt{SoszoroFej} hatás) validálása.}

\test{Sószóró fej működése - Sikertelen (Kifogyott só)}
{A sószóró fej áthalad az úton, de nincs benne só.}
{Ellenőrzi, hogy a sáv állapota változatlan marad és nem jön létre sózott réteg.}

\test{Zúzottkőszóró fej működése - Sikeres szórás}
{A hókotró zúzottkőszóró fejjel halad a jeges úton, és van készlete, amit kiszór.}
{Az új \texttt{ZuzottkoszoroFej} tesztelése: a fogyóeszköz csökken, és az adott sávon csökken a baleseti ráta (megcsúszás esélye).}

\test{Zúzottkőszóró fej működése - Sikertelen (Kifogyott kő)}
{A zúzottkőszóró fej halad, de üres a tartálya.}
{Validálja, hogy nem történik kőszórás, és a sáv csúszásveszélye nem csökken.}

\test{Söprő fej működése - Normál takarítás}
{A söprő fejjel haladó hókotró a csapadékot a menetirány szerinti jobbra lévő sávba tolja át.}
{A \texttt{SoproFej} és a \texttt{Sav} szomszédsági relációinak validálása: az aktuális sáv megtisztul, a szomszédoson nő a csapadék.}

\test{Jégtörő fej működése - Jégzúzás}
{A jégtörő fejjel felszerelt jármű szilárd jégpáncéllal borított sávon halad át.}
{Annak ellenőrzése, hogy a \texttt{JegtoroFej} a szilárd jeget feltört jéggé (törmelékké) alakítja, de fizikailag az úton hagyja.}

\test{Hányó fej működése - Teljes tisztítás}
{A nagyteljesítményű hányó fej áthalad egy csapadékos sávon.}
{Validálja, hogy a \texttt{HanyoFej} az úttól messzire repíti a havat, így az aktuális sáv teljesen mentesül a csapadéktól.}

% ==========================================
% 4. Környezeti hatások és játékállapot
% ==========================================

\test{Havazás szimulációja - Úton és hídon}
{Az \texttt{Uthalozat} globális havazást szimulál, amely érinti a nyílt útszakaszokat és a hidakat.}
{A folyamatos hógyarapodás ellenőrzése ezeken a \texttt{UtTipus} elemeken.}

\test{Havazás szimulációja - Alagútban}
{A globális havazás lefutása során ellenőrizzük az alagút típusú csomópontokat.}
{Validálja, hogy az alagút védett a csapadéktól, így ott nem nő a hóréteg vastagsága.}

\test{Leszórt só hatása (Olvadási idő működése)}
{Körök telnek el egy korábban már sózott útszakaszon.}
{A \texttt{Sav} állapotgépének tesztelése: a só megakadályozza a hó megmaradását, majd a rétegvastagsággal arányos idő után a só hatása elmúlik.}

\test{Játék vége - Kijárási tilalom (Vereség)}
{A folyamatos játék során a jeges utakon történő balesetek száma elér egy kritikus határértéket.}
{A vereséget kiváltó esemény (\texttt{Uthalozat} balesetszámláló) helyes detektálásának és a játék leállításának tesztelése.}

\test{Játék vége - Buszok elakadtak (Vereség)}
{A nagy hó miatt a város összes busza mozgásképtelenné (elakadttá) válik.}
{A játék végét hajtó alternatív logika ellenőrzése: a rendszer észleli a teljes forgalmi blokádot és lezárja a szimulációt.}

\section{Tesztelést támogató segéd- és fordítóprogramok specifikálása}
\comment{Specifikálni kell a tesztelést támogató segédprogramokat. Rövid bemutatással (elvárt funkcionalitás) specifikálni kell a tesztelést támogató segédprogramokat.}

A prototípus teszteléséhez nem használunk bonyolult, saját szintaktikájú pályaépítő vagy bemenet-generáló segédprogramokat, mivel a jelenlegi iterációban a konzolos interakciók és a soronkénti manuális bemenet-meghatározás elegendő a determinisztikus működéshez. A térkép inicializálását és a játékosok akcióit az előre megírt bemeneti parancsfájlok (input.txt) definiálják, amelyek a konzolos menüpontokra és kérdésekre adnak szekvenciális válaszokat.

A tesztelést és az értékelést az alábbi standard, illetve egyedi segédprogramok támogatják:

\subsection{Fordítóprogram}
A Java forráskódok lefordításához a standard Java Development Kit (JDK) \texttt{javac} fordítóprogramját használjuk. A fordítás parancssorból történik, amely a forrásfájlokból (\texttt{.java}) bájtkódot (\texttt{.class}) állít elő, fenntartva az eredeti csomagszerkezetet.

\subsection{Tesztfuttató szkript (\texttt{run\_test.bat} / \texttt{run\_test.sh})}
A tesztesetek automatizált és kényelmes futtatásáért egy platformfüggő parancsfájl (Windows alatt \texttt{.bat}, Linux/Unix alatt \texttt{.sh}) felel. 
\textbf{Funkcionalitás:}
\begin{itemize}
    \item Képes egyetlen, argumentumban megadott azonosítójú teszteset (pl. \texttt{run\_test.bat 4}), vagy az összes beépített teszt szekvenciális elindítására.
    \item A futtatás során a megfelelő Java parancsnak átadja az adott teszthez tartozó bemeneti fájlt (\texttt{--input=input.txt}).
    \item A program konzolos kimenetét teljes egészében egy szöveges fájlba (\texttt{--output=actual\_output.txt}) irányítja a későbbi kiértékelés céljából.
\end{itemize}

\subsection{Kimenet-ellenőrző (Diff) eszköz}
A tesztek sikerességének megállapításához nem írunk saját validáló programot, hanem az operációs rendszerek beépített szöveg-összehasonlító eszközeit alkalmazzuk (Windows környezetben az \texttt{fc}, Linux környezetben a \texttt{diff} parancs).
\textbf{Funkcionalitás:}
\begin{itemize}
    \item Összehasonlítja a tesztfuttató szkript által legenerált aktuális kimenetet (\texttt{actual\_output.txt}) a fejlesztők által manuálisan előállított és helyesnek ítélt elvárt kimenettel (\texttt{expected\_output.txt}).
    \item Eltérés esetén sorra pontosan megmutatja a különbségeket, míg egyezés esetén sikeres tesztlefutást igazol (0-s hibakód).
\end{itemize}
